\chapter{Grafos, digrafos y multigrafos}
\label{md:grafos}

\section{Vértices, aristas y variantes}

Salvo indicación contraria, se consideran grafos simples y finitos.
Un \textbf{grafo simple finito} $G=(V,E)$ consta de un conjunto finito
de vértices $V$ y un conjunto de aristas $E$, cada una de las cuales
es una pareja no ordenada de vértices distintos.
Los extremos de una arista son \textbf{adyacentes} y la arista es
\textbf{incidente} con ambos. El dibujo representa estas relaciones;
la posición o longitud dibujada de las aristas no forma parte del grafo.

En un \textbf{multigrafo} pueden existir varias aristas distintas
con los mismos extremos, llamadas paralelas. Hay que conservar
su identidad aunque conecten los mismos vértices.
Un \textbf{pseudografo} admite además lazos, que tienen ambos extremos
en el mismo vértice. Estas variantes pueden combinarse.

En un \textbf{digrafo} las aristas son dirigidas, o arcos:
$(u,v)$ va de $u$ a $v$, y no equivale a $(v,u)$.
Orientar un grafo consiste en elegir un sentido para cada arista.
Al recorrer un digrafo hay que respetar los sentidos de los arcos.

\section{Isomorfismo, grados y subgrafos}
\label{md:grados}

Dos grafos simples son \textbf{isomorfos} si existe una biyección entre
sus vértices que conserve la adyacencia en ambas direcciones.
Son la misma estructura con nombres de vértices diferentes.
El número de vértices, el de aristas y los grados se conservan:
una diferencia en esos datos descarta un isomorfismo, pero su igualdad
no basta para demostrarlo.
En las variantes también deben conservarse las multiplicidades,
los lazos y, en su caso, los sentidos.

El \textbf{grado} $\gr(v)$ de un vértice de un grafo simple es
el número de aristas incidentes con él. Un vértice de grado cero
es aislado. El recuento de extremos da el \textbf{lema del apretón de manos}:
\[
 \boxed{\sum_{v\in V}\gr(v)=2|E|.}
\]
Por tanto, el número de vértices de grado impar es par.
Si se cuentan grados en un pseudografo, un lazo contribuye dos veces,
una por cada extremo, y se mantiene esta identidad.

Un grafo es \textbf{$k$-regular} si todos los vértices tienen grado $k$;
entonces $k|V|=2|E|$.
El \textbf{grafo completo} $K_n$ une cada pareja de vértices distintos:
tiene grado $n-1$ en cada vértice y $n(n-1)/2$ aristas.

Un \textbf{subgrafo} se obtiene tomando algunos vértices y algunas
aristas cuyos extremos estén entre ellos. Al borrar un vértice se borran
sus aristas incidentes. El grado de un vértice en un subgrafo no supera
su grado original.

\section{Autómatas finitos}

Un \textbf{autómata finito determinista} se describe por
$\mathcal M=(Q,S,\delta,q_0,F)$:
\begin{itemize}
 \item $Q$ es el conjunto finito de estados y $S$ el alfabeto finito.
 \item $q_0\in Q$ es el estado inicial y $F\subseteq Q$ los estados finales.
 \item $\delta:Q\times S\to Q$ determina el siguiente estado al leer un símbolo.
\end{itemize}
Para procesar una palabra se parte de $q_0$ y se aplica sucesivamente
$\delta$ a cada símbolo. Se acepta la palabra si el estado alcanzado
pertenece a $F$; para la palabra vacía se permanece en $q_0$.

Su representación es un digrafo con estados como vértices y un arco
etiquetado $s$ de $q$ a $\delta(q,s)$ por cada símbolo.
Puede haber lazos o varios arcos con los mismos extremos.
El recorrido expresa la evolución del estado, no una elección arbitraria:
el símbolo leído fija la transición.
